\documentclass[showgrid=true]{swisstex}

\setrunningtitle{SwissTeX \textbar{} Minimalbeispiel}
\hypersetup{pdftitle={SwissTeX Minimalbeispiel}}

\begin{document}

\swisstitle{Minimalbeispiel}
  {SwissTeX\\Raster und Grundelemente}
  {Prüffall für Satzspiegel, Basislinienraster und Umbruchsteuerung}

\marg{Die Rasterkontrolle ist mit \texttt{showgrid=true} eingeschaltet}
\begin{lead}
Dieser Text steht im Vorspann. Er läuft auf 10,5\,pt über dieselbe
Rasterzeile von 13,5\,pt wie der Grundtext und bleibt deshalb in der Flucht.
\end{lead}

\section{Abschnitt mit Haarlinie}

\subsection{Grundtext}

Der Grundtext läuft im Flattersatz mit Trennung und optischem Randausgleich.
Absätze sind durch genau eine Rasterzeile getrennt, nicht durch einen Einzug.

Ein zweiter Absatz zeigt den Abstand. Alle vertikalen Masse sind ganzzahlige
Vielfache des Rastermasses, sodass die Grundlinien über Spalten und Seiten
hinweg fluchten.

\subsection{Formelsatz}

Formeln stehen linksbündig auf der Textachse:

\begin{gridblock}
\begin{equation*}
\gamma=\frac{\mathrm{d}n/\mathrm{d}T}{n-1}-\alpha .
\end{equation*}
\end{gridblock}

\section{Tabellen}

\subsection{Legende in der Marginalspalte}

\begin{swisstable}{Beispieltabelle mit Legende in der Marginalspalte}
\begin{tabularx}{\textcolumn}[t]{@{}S{28mm}@{\hspace{4mm}}L@{}}
\toprule
\tabhead{Merkmal} & \tabhead{Beschreibung} \\
\midrule
Satzspiegel & Der Tabellenkörper steht auf der Textachse und ist so breit
  wie der Grundtext. \\
\addlinespace[3.5pt]
Marginalspalte & Trägt Nummer und Titel der Tabelle, nie Satzmaterial. \\
\addlinespace[3.5pt]
Linien & Nur waagerechte Haarlinien, keine senkrechten. \\
\bottomrule
\end{tabularx}
\end{swisstable}

Nach der Tabelle setzt der Grundtext wieder auf dem Raster auf.

\colophon{Gesetzt mit SwissTeX. Grundschrift TeX Gyre Heros,
9,5/13,5\,pt, Satzspiegel 105\,mm, Marginalspalte 30\,mm.}

\end{document}
